\documentclass[11pt]{article}
\usepackage[margin=1in]{geometry}
\usepackage{amsmath,amssymb}
\usepackage{times}
\usepackage{hyperref}
\hypersetup{colorlinks=true,linkcolor=blue,urlcolor=blue,citecolor=blue}
\usepackage[numbers,sort&compress]{natbib}
\usepackage{enumitem}
\usepackage{titlesec}
\titleformat{\section}{\normalfont\large\bfseries}{\thesection}{1em}{}
\titleformat{\subsection}{\normalfont\normalsize\bfseries}{\thesubsection}{1em}{}

\title{\textbf{Counterfactual Before Classification: Vaccine Injury and the False Partition of Beneficiaries and Victims}}
\author{Flyxion \\ Independent Researcher}
\date{September 2026}

\begin{document}
\maketitle

\begin{abstract}
Public debate over vaccine safety is routinely conducted through a single rhetorical device: a ratio between an estimate of lives saved and an estimate of injuries caused. Defenders of vaccination use the ratio to argue that safety concerns are statistically negligible; critics use isolated injury counts to suggest the opposite. Both sides accept a shared and largely unexamined gloss on that ratio: that the vaccinated population divides cleanly into a set of beneficiaries and a set of victims, so that a favorable ratio amounts to a moral census of mostly-helped and rarely-harmed persons. That gloss is false, though the ratio itself is not. A serious adverse event and a prevented death need not belong to disjoint groups of people. They are potential components of the same person's counterfactual outcome, yet population-level marginal estimates do not reveal how often they coincide or what any particular injured person's outcome would have been without vaccination. This essay develops the argument in three stages: first, that ``vaccine injury'' equivocates between an individual causal claim and a population-level policy evaluation; second, that headline numerator figures for both harm and benefit are frequently miscombined across incompatible time windows, causal standards, and evidentiary sources; and third, and most fundamentally, that marginal benefit and harm counts, while legitimate for population-level policy comparison, cannot identify a partition between beneficiaries and victims, the extent of overlap between those categories, or the net causal outcome of any particular person classified as vaccine-injured.
\end{abstract}

\section{Two Scales of Description}

The phrase ``vaccine injury'' names two different kinds of claim depending on the scale at which it is asserted. At the individual scale, it is an ordinary causal proposition: this administration of this vaccine produced this adverse outcome in this person. Pharmacovigilance, clinical case review, and compensation adjudication all operate at this scale, and nothing about population-level reasoning invalidates it. A person can be genuinely injured by a vaccine in the same sense that a person can be genuinely injured by a medication, a surgical procedure, or a fall.

A vaccine is not, in the first instance, defined by its net population effect; the term names a biological intervention intended to induce protective immunity, and an ineffective or even net-harmful vaccine would remain a vaccine in that biological sense. But vaccination as a policy is evaluated and justified by its epidemiological function: the net change it produces in the distribution of disease outcomes across a population when deployed at scale. When a critic asserts that ``vaccine injury'' undermines the case for vaccination as a policy, the claim being made is necessarily a population-level claim, comparing the vaccinated world to the counterfactual unvaccinated world. But the evidence marshaled for that claim, an individual case report, a compensation claim total, a passive-surveillance count, is evidence generated at the individual scale. The rhetorical move consists of importing an individual-scale count into a population-scale argument without translating it, in the way an epidemiologist would, into a rate defined over a denominator of exposed persons, adjusted for detection intensity and compared against a background rate.

This is not merely a scale-mismatch of magnitude. It is closer to a category error, of the same kind as naming a firebreak a forest-destruction technology because building it requires cutting down some trees. The local description, trees were cut down, is materially true. The characterization is nonetheless taxonomically perverse, because it names the intervention by its comparatively small local cost while suppressing the much larger destruction the intervention exists to prevent. If a given vaccine prevents on the order of a thousand serious outcomes for every one it causes, describing the vaccine's social function as injurious, rather than as protective with a small and identifiable adverse tail, inverts the sign of the intervention it purports to characterize. The analogy is imperfect, since a firebreak's local cost is not itself a harm to the forest in the way a vaccine-attributable injury is a harm to a person; the point is only that naming an intervention by its comparatively small local cost, while suppressing the larger effect it exists to produce, misdescribes what the intervention is for.

None of this licenses dismissing individual harm as illusory. What it licenses is separating two questions that are routinely fused: whether this vaccine caused this outcome in this person, a question for clinical causal inference; and whether vaccination as a population intervention produces a net positive or net negative change in the distribution of serious outcomes, a question for population epidemiology. Conflating the two produces an argument that sounds quantitative while comparing quantities that are not commensurable.

\section{Surveillance, Compensation, and Causation}

A second source of confusion in the ``vaccine injury'' framing is that public commentary regularly treats several distinct kinds of data as interchangeable evidence of established harm. At least four are commonly run together: an adverse event occurring after vaccination; a passive-surveillance report of that event; a compensation claim; and a formally adjudicated causal attribution. Throughout, it also helps to keep two further senses separate: an adverse event is an outcome following vaccination that is causally attributable to it, whether serious or not; a net harm is the further claim that the person's counterfactual outcome was worse than it would have been without vaccination. The first is often directly observable in the individual case; the second, as Section~6 develops, is not.

Passive surveillance systems such as the U.S. Vaccine Adverse Event Reporting System exist precisely because they do not require prior proof of causation; their function is hypothesis generation and signal detection, which is why anyone, clinician or member of the public, can file a report describing any event that followed vaccination, verified or not. A raw VAERS count is therefore not a count of vaccine injuries; it is a count of temporally proximate events that someone thought worth reporting, and its size is sensitive to reporting behavior, media attention, and litigation activity as much as to true incidence. Compensation programs add a further layer, since eligibility criteria can be administrative or precautionary rather than a strict finding of scientific causation; the U.S. National Vaccine Injury Compensation Program, for instance, compensates certain injuries that fall within a predefined Vaccine Injury Table on a more permissive standard than proof of causation to a scientific certainty \citep{nvicp2024}. Reporting the total dollar value or case count of a compensation program as equivalent to a scientifically established injury count silently substitutes a legal or administrative standard for a causal one.

Formal causal adjudication looks different, and it is available for COVID-19 vaccines specifically. In 2024 a National Academies of Sciences, Engineering, and Medicine committee reviewed the epidemiological, clinical, and biological evidence for a defined list of candidate adverse events and reached specific, differentiated conclusions rather than a single verdict about ``vaccine injury'' in general. \citet{nasem2024} concluded that the evidence establishes a causal relationship between the mRNA vaccines, BNT162b2 and mRNA-1273, and myocarditis, while for numerous other candidate harms, including Guillain-Barr\'e syndrome, Bell's palsy, and, for the mRNA vaccines, thrombosis with thrombocytopenia syndrome, the committee found the evidence favored rejecting a causal relationship or judged it inadequate to accept or reject one. Notably, the committee did not rely solely on passive-surveillance totals; it treated an analysis of VAERS myocarditis cases, in which cases were validated by medical record review against CDC clinical criteria, as corroborating evidence alongside epidemiological cohort studies, precisely because raw report counts require independent verification before they support a causal conclusion \citep{nasem2024}.

This is what causal analysis of vaccine safety actually looks like: differentiated by product, by specific outcome, by evidentiary standard, and explicit about which associations are established, which are rejected, and which remain genuinely uncertain. A single undifferentiated figure for ``vaccine injuries,'' whatever its source, discards exactly the information that distinguishes a scientifically adequate safety claim from an inadequate one.

None of this makes ``vaccine injury'' an unscientific category as such. Tied to a specified product, outcome, causal standard, and exposed population, as in the National Academies conclusions above, it is a perfectly usable clinical and legal category, the ordinary object of pharmacovigilance and compensation adjudication. What is scientifically defective is a different move: treating that category as an undifferentiated epidemiological magnitude, or as a complete description of the net causal outcome for the person to whom it is attached, when in fact a documented adverse event is compatible with a range of counterfactual outcomes for that same person, as Section~6 develops.

\section{Observation Is Not Attribution}

An observed post-vaccination event, a clinically validated diagnosis, and a vaccine-attributable injury are not three names for the same datum. They are successive inferential objects. The first establishes temporal order: an event occurred after vaccination. The second establishes that the reported event satisfies the diagnostic criteria for a specified condition. The third establishes, at some stated level of confidence, that vaccination altered the probability of that condition relative to the relevant counterfactual.

Each transition can introduce classification error, and the two transitions are governed by different error structures rather than one continuous smoothing of the same uncertainty \citep{gustafson2003}. Passive-surveillance reports may include incomplete, duplicated, miscoded, or coincidental events. Clinical validation can reduce diagnostic misclassification without resolving causation. Conversely, an epidemiological increase in risk can support causal attribution at the population level without determining which particular cases would not have occurred in the absence of vaccination. The scientifically relevant numerator is therefore not obtained merely by counting reports and then assigning the resulting total a causal label.

Background incidence makes this distinction unavoidable. In a population containing millions of recently vaccinated people, some myocardial infarctions, strokes, neurological conditions, miscarriages, and deaths will occur shortly after vaccination even if vaccination has no causal effect on their frequency. The expected number of coincidental events may be large because the exposed population is large. Temporal proximity becomes informative only when combined with evidence such as an excess over the expected background rate, a reproducible risk window, dose or product specificity, biological plausibility, and the exclusion of important alternative explanations.

The movement from report to diagnosis and from diagnosis to causal attribution must consequently remain visible. Collapsing those stages into the single term ``vaccine injury'' conceals both the evidentiary work already accomplished and the uncertainty that remains.

\section{The Base-Rate Problem}

Safety signals are interpreted against an expected background distribution, not against an imaginary baseline in which no illness occurs after vaccination. Let $A$ denote an adverse event observed within a specified period after vaccination and let $V$ denote vaccination. The quantity
\[
P(A \mid V)
\]
is not, by itself, the probability that vaccination caused the event. Even a large number of observed events may provide little evidence of causation when the event is common in the population, the vaccinated population is enormous, or the observation window is sufficiently long.

The causal question concerns a contrast such as
\[
P(A \mid \text{do}(V=1)) - P(A \mid \text{do}(V=0)),
\]
not the mere frequency of $A$ among vaccinated persons \citep{pearl2009,pearl2010}. In observational data, estimating this contrast requires attention to age, health status, infection history, calendar time, health-care use, and other variables that affect both vaccination and event detection.\footnote{A fully Bayesian treatment of that estimation additionally propagates uncertainty in the background rate itself, rather than treating it as a fixed known quantity \citep{broemeling2013}; this affects how the contrast is estimated, not the definition of the contrast itself given in the text.} A surveillance count supplies neither the counterfactual probability nor the adjustment required to estimate it.

The same principle applies in the opposite direction. The fact that an event has a substantial background rate does not prove that every post-vaccination case is coincidental. A sufficiently concentrated risk window, an excess above baseline, product specificity, or convergent clinical and mechanistic evidence can shift the causal assessment. Base rates discipline inference; they do not predetermine it.

Public commentary commonly omits this inferential architecture. It presents the number of reports as though the relevant alternative were zero reports rather than the number expected among comparable unvaccinated people. The resulting figure may be numerically accurate as a report count while being scientifically uninformative as an injury estimate.

A closely related point concerns the operating characteristics of the surveillance system itself, not merely the background rate of the event it is trying to detect. A passive-reporting system can be thought of, loosely, as an imperfect diagnostic test for the presence of a true vaccine-caused case among the reports it receives, and the same logic that governs test accuracy governs it: sensitivity and specificity matter, but so does the prevalence of true cases among all the reasons a report might be filed \citep{broemeling2012}. Even a system built to capture nearly every genuine injury can have a low positive predictive value when true vaccine-attributable cases are rare relative to coincidental, misattributed, or duplicate reports, so that most individual reports, considered in isolation, are more likely to be noise than signal even while the system is doing exactly what it was designed to do. A report count is therefore not merely ``unverified'' in some generic sense; its evidentiary weight depends jointly on the detection system's operating characteristics and on the base rate of genuine causation, and neither can be read off the count alone.

\section{Decision Is Not Detection}

A further collapse compounds the ones already discussed: public debate routinely treats ``is vaccination safe'' as a single question with a single answer, when it is at least four distinct questions, each belonging to a different discipline, each answerable independently of the others.

Pharmacovigilance asks whether a causal signal exists at all: whether some vaccine, at some dose, produces some elevation in some outcome above what background rates and detection intensity would predict. Causal inference, given that a signal is credible, asks how large the effect is, in whom, and with what uncertainty, the question this essay takes up at length in the sections that follow. Risk analysis then asks how that estimated effect compares with the other effects, costs, and consequences at stake, weighing frequency against severity, certainty against uncertainty, and the vaccine's harms against the disease's harms and against the harms of available alternatives. Decision theory, finally, asks what action is preferable given all of the above together with the values and utilities attached to the possible outcomes, which is a further and irreducibly normative step rather than a further empirical one \citep{smith2010}.

These four questions are not stages of increasing certainty about a single fact; they are different kinds of question, and a negative or affirmative answer to one does not settle the others. Establishing that a causal signal exists, the pharmacovigilance question, does not by itself say whether the resulting risk is large enough to matter, the causal-inference question, whether it is large enough to change a policy recommendation once weighed against averted disease burden, the risk-analysis question, or what a decision-maker should therefore do, the decision-theory question. Conversely, a policy recommendation to continue a vaccination programme is not a denial that the pharmacovigilance question has a real, affirmative answer for some outcomes; a decision under uncertainty is not made by pretending the uncertainty, or the acknowledged cost, is zero.

Public commentary collapses all four questions into a single word, safety, and this collapse produces two mirror-image errors. The first treats any detected harm as automatically decisive against the policy, as though causal-inference or decision-theory work were unnecessary once a signal is confirmed. The second treats a favorable policy decision as evidence against the existence of any harm, as though the decision could only be correct if the pharmacovigilance question had come back negative. Both errors mistake a single stage of a four-stage inferential process for the whole of it, and the false partition this essay has been examining, in which every vaccinated person is either straightforwardly a beneficiary or straightforwardly a victim, is one particular symptom of that larger collapse: it is what results when the risk-analysis and decision-theory stages, which are precisely where heterogeneity, severity weighting, and population tradeoffs are supposed to be handled explicitly, are instead skipped and their conclusions read directly off the pharmacovigilance stage.

\section{Miscombined Numerators}

A further and more mundane problem afflicts the ratio argument even when it is deployed in vaccination's defense. Consider the widely cited estimate that COVID-19 vaccines prevented approximately 14.4 million deaths, rising to 19.8 million under an excess-mortality method, during the first year of the global vaccination programme \citep{watson2022}. That figure describes deaths averted across 185 countries and territories over a single specified twelve-month window, December 2020 to December 2021, using a transmission model fit against either officially reported COVID-19 deaths or excess-mortality estimates \citep{watson2022}. It is not a perpetual annual rate, and treating it as though vaccination saves ``14 million lives a year'' in every subsequent year overstates the claim the underlying study supports; the estimate is specific to a particular pandemic phase, variant mix, and baseline population immunity that do not recur unchanged.

Pairing that figure against a separately sourced injury count compounds the error, because the two numerators are rarely drawn from compatible populations, time windows, or causal standards. A benefit estimate derived from a formal transmission model fit to mortality data is not commensurable, without further adjustment, with a harm estimate drawn from a passive-reporting system, a compensation program, or a differently scoped time period. Dividing one headline number by the other, as in the intuitive calculation that a thousand lives are saved for every reported injury, produces an arithmetically correct ratio between two numbers that were never designed to be compared, in the same way that dividing a company's annual revenue by a competitor's quarterly revenue yields a number without yielding a meaningful comparison.

None of this weakens the underlying case for vaccination, which does not depend on this ratio; a modeled reduction of tens of millions of deaths against a background of specific, low-incidence, well-characterized adverse events remains an overwhelmingly favorable safety and efficacy profile under any reasonable accounting. The point is narrower: the popular one-line ratio is not itself a piece of evidence. It is a compressed and frequently mismatched pairing of two numbers, and its persuasive force trades on a precision the underlying data do not, in that specific combined form, possess.

\section{The False Partition}

The deeper problem, and the one that survives even after the individual-versus-population confusion and the miscombined-numerator problem are corrected, does not lie in the arithmetic of the ratio. A benefit-harm ratio does not, as a matter of mathematics, require that the saved and the injured be disjoint sets of people; a marginal comparison of expected counts remains well defined however much the underlying individuals overlap. The problem lies in a further step commonly taken with the ratio: reading ``a thousand people saved for every one injured'' as a moral census, a claim that mostly-distinct persons were helped and a small remainder were harmed. That reading is what fails, and it fails for a structural reason rather than an incidental one.

Let $Y_i(1)$ denote person $i$'s outcome under vaccination and $Y_i(0)$ denote the same person's outcome had they remained unvaccinated. The individual causal effect of vaccination on person $i$ is the counterfactual contrast
\[
\Delta_i = Y_i(1) - Y_i(0),
\]
and only one of the two potential outcomes, $Y_i(1)$ or $Y_i(0)$, is ever observed for any given person; the other is, in Holland's phrase, the fundamental problem of causal inference \citep{holland1986}. Population studies estimate the missing counterfactual on average across many people; they do not, and structurally cannot, report it person by person.

A person who develops vaccine-attributable myocarditis after an mRNA vaccine has experienced a genuine, causally attributable adverse event: $Y_i(1)$ includes myocarditis, and this component can be clinically documented in the individual. That fact alone does not fix $Y_i(0)$. If that same person would have contracted a severe or fatal case of COVID-19 without vaccination, $Y_i(0)$ would have been worse than $Y_i(1)$ despite the adverse event, and the net counterfactual effect $\Delta_i$ for that person would still be favorable. But $Y_i(0)$ is not observed; the protective component, unlike the adverse component, remains a counterfactual possibility whose frequency can be estimated only at the population level, never verified in the individual case. Classification as vaccine-injured therefore does not establish that a person was made worse off overall, but neither does population-level vaccine effectiveness establish that any particular injured person was among those saved. The potential outcomes represented by ``injured'' and ``saved'' are not logically exclusive; someone who experiences a vaccine-attributable adverse event may also be someone for whom vaccination prevented a more serious outcome. Because the latter claim concerns an unobserved counterfactual, however, the overlap generally cannot be identified person by person and should not be asserted as though it were an observed classification (the adverse-event / net-harm distinction of Section~2).

There is a condition under which the partition would be licensed: monotonicity, the assumption that vaccination never makes any person worse off on the outcome in question, so that $Y_i(1) \geq Y_i(0)$ for every $i$. Under monotonicity, a person who experiences an adverse event cannot also be someone for whom vaccination averted a worse outcome on that same dimension, because no one is made worse off by exposure; the ``injured'' and ``saved'' categories would then be disjoint by construction, and the ratio could be read as something closer to a partition. The difficulty is that monotonicity is a substantive assumption about effect heterogeneity, not a mathematical truth, and it is exactly what is at issue for vaccines across heterogeneous outcomes. The National Academies findings reviewed in Section~2 already indicate that it does not hold uniformly: a causal relationship is established for the mRNA vaccines and myocarditis \citep{nasem2024}, which is to say that for some persons $Y_i(1)$ is worse than $Y_i(0)$ on that dimension. Whether the same persons are those for whom vaccination averted a worse outcome on a different dimension is a further question, and one the monotonicity assumption would need to settle in order to license the partition. Absent that assumption, the disjointness of beneficiaries and victims is not a finding but a premise, and one the data do not supply. The false partition, in other words, is precisely what monotonicity would buy; the substantive claim of this essay is that no warrant for it exists here.

At the population level, writing a person's net outcome as $U_i = B_i - H_i$, with $B_i$ the benefit and $H_i$ the harm expressed in commensurable units,
\[
\mathbb{E}[U_i] = \mathbb{E}[B_i] - \mathbb{E}[H_i]
\]
by linearity of expectation, and this identity holds whether or not $B_i$ and $H_i$ are independent or concentrated in the same individuals. Overlap does not invalidate this subtraction; it invalidates reading the subtraction as a partition of persons. Nor does the joint distribution of $Y_i(1)$ and $Y_i(0)$ become recoverable through some more refined statistical technique. This is not a limitation of sample size or estimator quality: the two potential outcomes are never jointly realized in the same unit, so the joint distribution is unidentified without additional structural assumptions, monotonicity among them, along with a model of effect heterogeneity or a valid instrument, none of which is generally defensible for vaccine outcomes across heterogeneous harms.

\section{Events Are Not Persons}

Much of the apparent paradox arises because public discussion moves without warning between counting events and classifying persons. A benefit-harm analysis may estimate deaths averted, hospital admissions prevented, cases of myocarditis caused, and episodes of anaphylaxis caused. These are counts of outcomes, not necessarily counts of mutually exclusive human beings. One person may contribute more than one outcome to such an accounting, just as the same patient may experience an adverse drug reaction while also receiving a therapeutic benefit from the drug.

Consequently, a statement such as ``one serious adverse event occurred for every thousand deaths prevented'' does not logically imply that one person was sacrificed so that a thousand different people could live. Properly interpreted, it compares the estimated frequencies of two kinds of causal event. Its common paraphrase, ``a thousand people were saved for every person injured,'' introduces an additional claim: that the corresponding people form distinct classes. That claim is neither required by the event ratio nor generally established by the data used to calculate it.

The distinction matters because events and persons answer different questions. Counting causally attributable events helps determine whether deploying an intervention improves population health. Classifying a person as a net beneficiary or net victim would require comparing that person's entire outcome under vaccination with the outcome that same person would have experienced without vaccination. The former can often be estimated from population data. The latter ordinarily cannot be observed.

Thus, overlap does not make aggregate comparison impossible. It makes a particular interpretation of aggregate comparison invalid. The error lies not in counting benefits and harms but in converting marginal event counts into a partition of the population.

\section{Subtraction Does Not Require Disjointness}

The identity established in Section~6 explains why aggregate subtraction survives overlap between benefit and harm. What it leaves open is the joint distribution connecting the two. Two vaccination programmes could produce the same aggregate net benefit while distributing their effects very differently. In one programme, the people experiencing adverse effects might also be those receiving the greatest protection from disease. In another, adverse effects might be concentrated among people at very low risk from the disease and therefore unlikely to receive much direct benefit. The same marginal totals could conceal substantially different distributions of individual welfare.

The aggregate expression
\[
N_{\text{harm prevented}} - N_{\text{harm caused}}
\]
is therefore not mathematically defective merely because some people may appear substantively on both sides. Its limitation is interpretive. Unless the harms being prevented and caused have been placed on a common severity scale, the subtraction may combine incommensurable outcomes. Even when commensurability has been supplied through quality-adjusted life-years, disability-adjusted life-years, expected mortality, or another utility measure, the result describes average population consequences rather than identifying the net causal outcome of each person.

The failure of the familiar ratio is consequently more precise than simple invalidity. A properly constructed event ratio can summarize population-level tradeoffs. It cannot disclose whether its benefits and harms occur in the same people, whether a particular injured person was nevertheless a net beneficiary, or whether the intervention distributes its small adverse tail equitably.

\section{The Joint Distribution Is Partially Unidentified}

The preceding point can be given an exact form. Let $B_i$ indicate that vaccination prevents a serious outcome for person $i$, and let $H_i$ indicate that vaccination causes a serious adverse event in that person. Population studies may estimate the marginal probabilities
\[
P(B_i = 1) \qquad \text{and} \qquad P(H_i = 1),
\]
yet those quantities do not by themselves determine the joint probability
\[
P(B_i = 1, H_i = 1).
\]
Many different joint distributions are compatible with the same two marginal probabilities \citep{morganwinship2015}. The benefits and harms might occur almost entirely in different people, overlap extensively, or exhibit an intermediate structure without changing either headline count.

The possible amount of overlap is bounded by the Fr\'echet inequalities:
\[
\max\{0, \, P(B_i=1) + P(H_i=1) - 1\} \; \leq \; P(B_i=1, H_i=1) \; \leq \; \min\{P(B_i=1), P(H_i=1)\}.
\]
These bounds reveal what the marginal counts leave unresolved. If a benefit occurs in ten per cent of recipients and an adverse effect in one per cent, the proportion experiencing both may lie anywhere between zero and one per cent unless further assumptions or evidence constrain the dependence between them. It is worth being explicit that the bounds cut both ways. The lower bound, zero overlap, is the case the naive partition assumes; the upper bound, complete overlap of the smaller category within the larger, is the case most favorable to this essay's suggestion that injured persons may often also be net beneficiaries. The marginal counts alone do not favor either end over the other, or anything in between; the bounds establish only that the data available to the public debate cannot settle where in that range the truth lies, not that the truth is likely to sit anywhere in particular within it.

The difficulty is even greater because $B_i$ is ordinarily counterfactual. Whether vaccination prevented severe disease in a given recipient depends on the unobserved outcome that would have followed non-vaccination. Consequently, the joint distribution cannot usually be recovered merely by recording more outcomes in vaccinated recipients. Its estimation requires assumptions about causal structure, exchangeability, exposure, infection risk, competing events, and the relation between susceptibility to disease and susceptibility to adverse effects \citep{pearl2009}.

This underidentification does not prevent estimation of average population benefit. Nor does it invalidate subtraction of expected benefits and harms, as Section~8 shows. It prevents the marginal totals from being interpreted as a complete map of who benefited, who was harmed, and who experienced both. The familiar ratio compresses two marginal estimates while leaving their dependence structure, and therefore the distribution of individual net outcomes, unspecified.

\section{The Unobservability of the Saved Individual}

The expression ``lives saved'' is itself easily misunderstood. Except in unusual clinical circumstances, an averted death is not an observable event attached to an identifiable person. Researchers observe mortality under one exposure condition and estimate how much mortality would have occurred under another. The difference between those population distributions is called deaths averted, but the analysis ordinarily cannot identify which particular living individuals would otherwise have died.

This creates an important asymmetry between the two sides of the public comparison, though not the simple asymmetry of observed harm against unobserved benefit. A case of vaccine-attributable myocarditis can sometimes be identified clinically in a particular person through its timing, presentation, exclusion of alternatives, and consistency with established epidemiological evidence \citep{nasem2024}. That much is directly attestable. But the adverse event's status as a net harm is not: whether the person's overall counterfactual outcome was worse with vaccination than without it remains, in this respect, exactly as unobservable as whether vaccination saved them. Both net effects are counterfactual. What is asymmetric is that the adverse event can be documented in the individual case while the offsetting protective effect cannot, so only one direction of the event is directly attestable even though neither direction of the net effect is.

It would therefore be excessive to say that a person with a vaccine-attributable injury was also demonstrably saved by the vaccine. The defensible statement is conditional: that person might have experienced a still worse outcome without vaccination, and classification by observed adverse effect alone does not determine the sign of the person's total causal outcome. Conversely, the vaccine's large average protective effect does not prove that every person who experienced an adverse effect was personally better off after receiving it.

This is an instance of the fundamental problem of causal inference rather than an evidentiary defect peculiar to vaccines \citep{holland1986}. For each person, either the vaccinated or unvaccinated trajectory becomes observable, never both. Randomized trials and observational causal methods can estimate differences between distributions under suitable assumptions, but they do not reconstruct both histories for the same individual. The population may therefore contain people who were protected without injury, injured while still receiving net protection, injured without receiving an offsetting benefit, and neither injured nor materially protected. The proportions occupying those latent categories are not automatically recoverable from the two marginal headline counts.

\section{Population Benefit Is Not Universal Individual Benefit}

A strongly favorable population result does not imply that vaccination benefited every recipient. Population averages permit heterogeneity. Some recipients receive substantial protection, some receive little incremental protection because they would not otherwise have become infected or seriously ill, and a very small number experience adverse outcomes that may exceed whatever direct benefit they received. Recognizing the existence of the final group does not contradict the conclusion that vaccination is strongly beneficial at population scale.

The converse is equally important. Establishing that a particular person suffered a genuine vaccine-attributable injury does not establish that the vaccination programme was harmful, that the vaccine's overall benefit-risk balance was unfavorable, or even that the injured person's complete counterfactual outcome was worse than it would have been without vaccination. These are different causal propositions requiring different evidence. An individual causal finding concerns one pathway from exposure to harm. A policy judgment concerns the distribution of all relevant pathways across the exposed and counterfactual populations.

This prevents two symmetrical errors. The first uses aggregate benefit to deny individual injury: because vaccination saved millions of lives, a particular claimant could not have been harmed. That inference is false because favorable averages permit rare negative effects. The second uses individual injury to reverse the population verdict: because this claimant was harmed, vaccination was therefore injurious as a social intervention. That inference is equally false because an individual adverse outcome does not determine the sign or magnitude of the aggregate causal effect.

The appropriate vocabulary must preserve both truths simultaneously. A vaccine may possess an overwhelmingly favorable population benefit-risk profile while causing a small number of serious individual adverse outcomes. There is no contradiction because the propositions operate at different levels. The contradiction is manufactured only when an individual adverse effect is treated as a complete characterization of population function, or when population success is treated as evidence that no individual loss occurred.

\section{Compensation Is Not a Verdict Against Vaccination}

Recognition and compensation of vaccine-attributable injury are sometimes presented as concessions that undermine confidence in vaccination. This interpretation mistakes the ethical allocation of residual risk for a judgment about the intervention's aggregate value. If society promotes an intervention because its distributed benefits greatly exceed its distributed harms, it may acquire a special obligation toward the small number of people who bear serious costs on behalf of that population policy.

A compensation programme can therefore express confidence in vaccination rather than doubt about it. Its underlying argument is that the intervention is sufficiently beneficial to warrant collective deployment, but that the uncommon costs of deployment should not remain concentrated arbitrarily in the individuals upon whom they happen to fall. Compensation socializes part of the adverse tail just as vaccination socializes protection against infection. It does not transform a favorable benefit-risk distribution into an unfavorable one.

Nor does compensation necessarily provide the same kind of causal evidence as an epidemiological study. Administrative systems may deliberately adopt standards that are more permissive than scientific proof beyond reasonable doubt, because their purpose is remedial rather than purely explanatory. The U.S. National Vaccine Injury Compensation Program illustrates the point directly: an injury listed on its Vaccine Injury Table within the specified onset window is presumed caused by the vaccine without further proof, while an injury not on the table must instead be shown by a preponderance of the evidence, a more likely than not standard well short of the causal certainty the National Academies review applies \citep{nvicp2024,nasem2024}. A compensated claim must therefore be described according to the rules under which it was accepted, not automatically converted into a universally established scientific case.

The ethically serious position is neither to erase injured individuals beneath population averages nor to let exceptional injuries define the population intervention. It is to preserve causal specificity, compare benefits and harms in compatible terms, reduce avoidable risks through product selection and dosing policy, and compensate severe residual harms where appropriate. Scientific recognition of rare injury and epidemiological recognition of overwhelming net benefit are not rival narratives. They are complementary descriptions of an intervention whose effects are heterogeneous but whose population-level sign can nevertheless be clear.

\section{Risk Stratification Does Not Reverse the Population Result}

A favorable population average may conceal meaningful variation by age, sex, previous infection, comorbidity, vaccine product, dose number, and epidemic phase. This heterogeneity is particularly relevant when the disease risk and the adverse-event risk are distributed differently. A benefit-risk estimate calculated across the entire population may not describe the balance for a young person at low risk of severe disease as accurately as it describes the balance for an older person with substantial clinical vulnerability.

Acknowledging such variation does not vindicate an undifferentiated category of vaccine injury. It requires greater differentiation, not less. The appropriate response is to estimate product- and subgroup-specific absolute risks, alter dose intervals or product recommendations where warranted, and reconsider recommendations as baseline immunity and disease incidence change. An intervention can remain strongly beneficial overall while requiring narrower targeting in particular circumstances.

This also explains why an estimate from the first pandemic year cannot function as a timeless constant \citep{watson2022}. The counterfactual changes as population immunity accumulates, variants change, treatments improve, and the age distribution of recipients shifts. Vaccine benefit is not an intrinsic number permanently attached to a vial; it is a causal contrast between vaccination and non-vaccination under specified conditions. Vaccine-attributable risk may also change with product, dose, schedule, and recipient characteristics. A scientifically meaningful comparison must therefore be indexed to a population, intervention, comparator, outcome, and time.

The demand for stratification should not be confused with the claim that aggregate evidence is useless. Population estimates answer population questions, while subgroup estimates refine decisions where heterogeneous risks make the aggregate insufficient. The conceptual mistake occurs when either scale is made to answer the other's question.

\section{What the Correct Comparison Requires}

Bringing the preceding distinctions together, a scientifically adequate comparison of vaccine benefit and vaccine harm requires several things that the popular framing routinely omits.

It requires a shared population and time window for both the benefit and harm estimates, rather than pairing a multi-country annual mortality model with an open-ended cumulative surveillance total. It requires the harm term to be built from causally adjudicated events, of the kind the National Academies review provides for specific vaccine-outcome pairs, rather than from unverified reports whose count reflects reporting behavior as much as true incidence \citep{nasem2024}. It requires the benefit term to be scoped to the period and population it actually models, rather than extrapolated as a recurring annual figure \citep{watson2022}. And it requires that whatever aggregate net figure results be reported as exactly that, a population-level marginal balance, not as a census identifying which persons were net beneficiaries, which were net victims, or how much the two categories overlap, since that further claim is not something the aggregate can support.

Severity asymmetry also needs to be kept explicit rather than collapsed into a single count. Death is irreversible in a way that most, though not all, vaccine-attributable adverse events are not; severity, duration, and prognosis vary considerably even among the small number of causally established harms, and myocarditis itself ranges from a mild, self-limited episode to a serious cardiac event, so this variation should be represented rather than assumed away in either direction. The National Academies review's differentiated, outcome-specific conclusions already reflect the kind of accounting that a single ``vaccine injury'' figure destroys \citep{nasem2024}; a responsible comparison keeps these terms separate rather than collapsing them into one undifferentiated count on each side of the ledger.

None of this is a demand for a more favorable conclusion. It is a demand for the conclusion, whatever it turns out to be for a given vaccine, outcome, and population, to be reached by comparing quantities that are actually comparable, and by not asking an aggregate marginal comparison to answer a question about individuals that it cannot answer. The available evidence for COVID-19 vaccines, properly scoped, supports a strongly favorable safety and efficacy profile at the population level; that conclusion does not depend on, and is not strengthened by, a ratio recast as a partition of the population into the saved and the injured.

\section{Present Harm and Absent Benefit}

The preceding sections explain why the popular comparison is statistically incomplete. They do not explain why the incompleteness so consistently favors the injury narrative over the benefit narrative in public perception, rather than failing in some more neutral or random way. Part of the answer lies not in statistics but in what each side of the comparison leaves behind.

A prevented event has no phenomenology. No one experiences the hospitalization that did not occur, notices the ventilator that was never required, or can name the date on which vaccination prevented their death. Prevention appears socially as an absence, and an absence does not spontaneously generate testimony. An adverse event has the opposite structure: it is experienced, dated, diagnosed, and can be narrated as a personal history. This asymmetry in what each outcome leaves behind holds even when the underlying benefit-risk balance is strongly favorable. The intervention's costs produce witnesses; its successes typically remove the very events that would have produced witnesses to its necessity.

The resulting public archive of vaccine stories is consequently not a neutral sample of the intervention's causal effects. It is selectively populated by the outcomes capable of becoming stories. A person who becomes ill following vaccination has an immediate reason to search for an explanation, seek out others with similar experiences, or come forward publicly. A person who remained healthy because vaccination prevented a severe infection typically has no comparable occasion to attribute anything to the vaccine at all, since nothing happened to prompt the attribution. The absence of comparable beneficiary testimony does not show that the benefits are less real than the harms; it shows only that prevention and injury have different conditions of visibility, independent of their relative frequency or magnitude. This visibility asymmetry is separate from, and compounds, the more general finding that losses tend to be weighted more heavily than commensurate gains in human judgment \citep{kahneman1979}: an outcome that is both a loss and the only side of the ledger with a witness is doubly, not merely additively, disadvantaged in public perception.

This is one reason restoring the exposure denominator, the causal standard, and the counterfactual comparison, as this essay has argued throughout, is more than a technical correction. It reconstructs a class of outcomes that direct experience, by its nature, cannot supply testimony for. Epidemiology is what makes prevented outcomes inferentially visible at all; where confidence in that kind of institutional inference is weak, the asymmetry described here gives the visible, narratable side of the comparison a standing advantage that has nothing to do with which side is larger. The asymmetry has an epidemiological counterpart in Rose's prevention paradox: a population-level preventive measure can produce substantial aggregate benefit while providing little perceptible benefit to most of the people who participate in it \citep{rose1981,rose1992}. Vaccination sharpens the epistemic version of that paradox, because the prevented outcome, unlike a merely small individual benefit, leaves no individual experience through which its prevention could be recognized at all.

\section{Epistemic Dependence and the Social Amplification of Risk}

The visibility asymmetry of the preceding section explains why injury and prevention leave different kinds of trace. It does not yet explain why an audience presented with both the trace and its statistical correction so often weights the trace more heavily. Two further, independently documented mechanisms bear on that question: the structure of epistemic dependence itself, and the social amplification of risk signals as they pass through public communication.

Most people cannot personally reproduce a vaccine trial, audit a multi-country mortality model, or estimate an adverse event's counterfactual background incidence. Knowledge at this scale is necessarily distributed, and rational belief on such matters consequently depends on warranted reliance on experts whose evidence one cannot personally possess or independently reconstruct \citep{hardwig1985}. This dependence is not a defect to be overcome; refusing it in favor of what one can verify unaided would leave a person with only the narrow, informal, and largely untested beliefs that fall within a single lifetime's direct experience.

That structural dependence interacts with a further asymmetry in how evidence is processed once received. Slovic and colleagues distinguish ``risk as feelings,'' organized through images, affect, and immediate association, from ``risk as analysis,'' organized through probabilities and formal comparison of the kind this essay has developed throughout \citep{slovic2004}. A personal injury account is well suited to the first mode: it arrives with a face, a chronology, and an affectively vivid narrative. A modeled reduction in mortality is well suited only to the second: it is a difference between counterfactual population distributions with no comparable affective handle. Where confidence in the institutions that produce analytic risk information is low, the feelings-based mode does not merely supplement the analytic one but can displace it in public judgment.

Media and communication institutions then act on this asymmetry rather than merely reporting around it. Kasperson and colleagues describe how risk signals are selected, interpreted, repeated, and transmitted through social and institutional ``amplification stations,'' with the result that a signal's public salience can diverge substantially from its statistical magnitude or causal certainty \citep{kasperson1988}. An isolated injury story, reported without its exposure denominator, background rate, or causal standard, is exactly the kind of signal this framework predicts will be amplified: it requires no institutional trust to find compelling, only a witness.

A closely related psychological finding, the identifiable-victim effect, is often invoked here and deserves a more careful statement than it usually receives. Jenni and Loewenstein's classic account holds that an individuated victim can elicit a stronger response than a statistically described group facing the same aggregate harm \citep{jenni1997,small2007}. The analogy to vaccine injury should not be overstated, however: the experimental literature concerns donation and helping behavior rather than causal attribution, and more recent, well-powered preregistered replications have failed to reproduce the original identifiable-victim findings, suggesting the effect is considerably less uniform than its popular presentation implies \citep{maier2023}. What survives that caveat is the narrower structural point rather than a specific effect size: injury is capable of individuation in a way that successful, population-level prevention structurally is not, regardless of how reliably that individuation translates into a measurable shift in judgment in any particular experimental setting.

None of this licenses a demand that the public simply defer more readily to institutional authority. Institutions partly earn distrust through opacity, inconsistency, exaggerated certainty, and failure to acknowledge error, and the correct response to that earned distrust is not louder assurance. O'Neill's distinction between trust and trustworthiness is the relevant one here: an institution's task is not chiefly to cultivate more trusting attitudes but to become assessably worthy of reliance, through intelligible evidence, accountability, and practices that a non-expert can use to judge competence and good faith even without verifying the underlying science directly \citep{oneill2002}. Restoring the denominator, the causal standard, and the exposure population, the corrective this essay has argued for throughout, is one instance of exactly that kind of assessable practice; it does not ask the public to trust the conclusion, only to be given what is needed to check the reasoning behind it.

\section{The Sociology of Visibility: From Testimonial Asymmetry to Nocebo Mechanisms}

Section~16 made an epistemological claim: prevention and injury are asymmetric in the testimony they generate, because prevention leaves no narratable personal experience while an adverse event can be experienced, diagnosed, and converted into public testimony. This section shifts to a different, empirically studied question that compounds rather than restates that asymmetry: not only does injury produce testimony where prevention does not, but the presence of that testimony can itself increase the rate at which further injury is reported, through the well-documented phenomenon of the socially transmitted nocebo effect.

The nocebo effect denotes adverse outcomes produced by negative expectation rather than by a treatment's pharmacological action. In COVID-19 vaccine trials specifically, a systematic review and meta-analysis of twelve randomized, placebo-controlled trials found that 35.2 percent of placebo recipients reported at least one systemic adverse event after the first dose and 31.8 percent after the second, headache and fatigue the most common; the same analysis estimated that roughly three-quarters of the systemic adverse events reported by vaccine recipients were attributable to this nocebo component rather than to the vaccine's pharmacological action \citep{haas2022}. Nocebo responses are not confined to classical conditioning or a clinician's verbal warning. A separate line of experimental and survey research finds that expectations of harm can also be transmitted through social observation, including exposure to other recipients' reported side effects on social media, and that this social-learning pathway measurably increases the likelihood that a subsequent vaccine recipient goes on to report side effects of their own \citep{tan2022}.

This mechanism has two consequences for the argument of this essay. First, it supplies a causal pathway, not merely a descriptive correlation, for why injury narratives carry disproportionate public weight: a widely circulated account of harm can, through ordinary expectancy effects, increase the rate at which subsequent recipients experience and report symptoms themselves, feeding a narrative-expectation-symptom cycle that has no benefit-side counterpart, since there is no comparable mechanism by which hearing about someone else's averted infection would make one's own vaccination more protective. Second, because the mechanism is measurable, the asymmetry it compounds is not simply an unavoidable fact about testimony; it is amplified or damped by how side effects are communicated in the first place, and expectation-sensitive communication has independently been studied as a lever on nocebo symptom reporting.

None of this is offered as an explanation of vaccine injury in general, and it should not be read as reclassifying the causally adjudicated harms discussed earlier, the mRNA-myocarditis association of Section~2 among them, as nocebo artifacts; those remain real, pharmacologically caused effects regardless of how much additional nocebo burden surrounds them. The point is narrower and complements rather than substitutes for the statistical argument of the rest of this essay: once genuine causal harms have been separated from expectation-driven ones by the kind of evidentiary work Sections~2 through~5 describe, how side effects are described in public communication becomes itself a variable that affects reported outcomes, not merely a channel that reports them. Leaving that variable unmanaged is not a neutral default; it is a choice to let the testimonial asymmetry of Section~16 operate unchecked.

\section{Conclusion}

Benefit-harm ratios are not meaningless merely because benefits and harms may overlap within persons. They are meaningful as comparisons of population-level marginal outcomes, provided their populations, time windows, causal standards, and units are compatible, and the case for COVID-19 vaccination does not need a disjoint-population reading to be strong. What such ratios cannot establish is a partition between beneficiaries and victims, the amount of overlap between those categories, or the net causal outcome of any particular person classified as vaccine-injured. ``Vaccine injury,'' used as a specified clinical and legal category tied to a product, an outcome, and a causal standard, is a legitimate scientific object, and formal causal review already identifies specific, real, and differentiated harms of that kind. What is scientifically defective is a further and usually implicit step: treating an aggregate net-benefit figure as though it resolved, for any given injured person, whether that person was on balance helped or harmed, when the joint distribution connecting a documented adverse event to an unobserved counterfactual is precisely what population data cannot recover.

The corrective this essay has argued for is therefore not only statistical. A comparison stated with every qualification defended above, restricted to compatible populations and time windows, built from causally adjudicated rather than merely reported harms, and honest about the non-identifiability of the joint distribution, still has to be received by an audience, and Sections~16 through~18 give three compounding reasons why that reception will not be neutral: prevention structurally produces no witnesses while injury does, that asymmetry is processed through modes of judgment and channels of transmission that favor vivid individual testimony over statistical analysis, and the resulting asymmetry can be further amplified, through ordinary expectancy mechanisms, by the very act of reporting on it. Correcting the ratio is necessary but not sufficient. The remaining work is rhetorical and institutional as much as statistical: reporting practices that make the corrected comparison, and not only the vivid half of it, the one an audience actually encounters. The corrective is not a different ratio. It is dropping the classification the ratio is asked, beyond its actual content, to support, and building the communication of both halves of the comparison so that dropping it is possible in practice and not only in principle.

\bibliographystyle{plainnat}
\begin{thebibliography}{9}

\bibitem[Broemeling(2012)]{broemeling2012}
Broemeling, L.~D. (2012). \textit{Advanced Bayesian Methods for Medical Test Accuracy}. Boca Raton, FL: CRC Press.

\bibitem[Broemeling(2013)]{broemeling2013}
Broemeling, L.~D. (2013). \textit{Bayesian Methods in Epidemiology}. Boca Raton, FL: Chapman \& Hall/CRC.

\bibitem[Gustafson(2003)]{gustafson2003}
Gustafson, P. (2003). \textit{Measurement Error and Misclassification in Statistics and Epidemiology: Impacts and Bayesian Adjustments}. Boca Raton, FL: Chapman \& Hall/CRC.

\bibitem[Haas et al.(2022)]{haas2022}
Haas, J.~W., Bender, F.~L., Ballou, S., Kelley, J.~M., Wilhelm, M., Miller, F.~G., Rief, W., \& Kaptchuk, T.~J. (2022). Frequency of adverse events in the placebo arms of COVID-19 vaccine trials: a systematic review and meta-analysis. \textit{JAMA Network Open}, 5(1), e2143955.

\bibitem[Hardwig(1985)]{hardwig1985}
Hardwig, J. (1985). Epistemic dependence. \textit{The Journal of Philosophy}, 82(7), 335--349.

\bibitem[Holland(1986)]{holland1986}
Holland, P.~W. (1986). Statistics and causal inference. \textit{Journal of the American Statistical Association}, 81(396), 945--960.

\bibitem[Jenni and Loewenstein(1997)]{jenni1997}
Jenni, K.~E., \& Loewenstein, G. (1997). Explaining the identifiable victim effect. \textit{Journal of Risk and Uncertainty}, 14(3), 235--257.

\bibitem[Kahneman and Tversky(1979)]{kahneman1979}
Kahneman, D., \& Tversky, A. (1979). Prospect theory: an analysis of decision under risk. \textit{Econometrica}, 47(2), 263--291.

\bibitem[Kasperson et al.(1988)]{kasperson1988}
Kasperson, R.~E., Renn, O., Slovic, P., Brown, H.~S., Emel, J., Goble, R., Kasperson, J.~X., \& Ratick, S. (1988). The social amplification of risk: a conceptual framework. \textit{Risk Analysis}, 8(2), 177--187.

\bibitem[Maier et al.(2023)]{maier2023}
Maier, M., Wong, Y.~C., \& Feldman, G. (2023). Revisiting and rethinking the identifiable victim effect: replication and extension of Small, Loewenstein, and Slovic (2007). \textit{Collabra: Psychology}, 9(1), 90203.

\bibitem[Morgan and Winship(2015)]{morganwinship2015}
Morgan, S.~L., \& Winship, C. (2015). \textit{Counterfactuals and Causal Inference: Methods and Principles for Social Research} (2nd ed.). Cambridge: Cambridge University Press.

\bibitem[Health Resources and Services Administration(2024)]{nvicp2024}
Health Resources and Services Administration. (2024). \textit{National Vaccine Injury Compensation Program}. U.S. Department of Health and Human Services. \url{https://www.hrsa.gov/vaccine-compensation}

\bibitem[National Academies of Sciences, Engineering, and Medicine(2024)]{nasem2024}
National Academies of Sciences, Engineering, and Medicine. (2024). \textit{Evidence Review of the Adverse Effects of COVID-19 Vaccination and Intramuscular Vaccine Administration}. Washington, DC: The National Academies Press. \url{https://doi.org/10.17226/27746}

\bibitem[O'Neill(2002)]{oneill2002}
O'Neill, O. (2002). \textit{A Question of Trust: The BBC Reith Lectures 2002}. Cambridge: Cambridge University Press.

\bibitem[Pearl(2009)]{pearl2009}
Pearl, J. (2009). \textit{Causality: Models, Reasoning, and Inference} (2nd ed.). Cambridge: Cambridge University Press.

\bibitem[Pearl(2010)]{pearl2010}
Pearl, J. (2010). An introduction to causal inference. \textit{The International Journal of Biostatistics}, 6(2), Article 7.

\bibitem[Rose(1981)]{rose1981}
Rose, G. (1981). Strategy of prevention: lessons from cardiovascular disease. \textit{British Medical Journal}, 282, 1847--1851.

\bibitem[Rose(1992)]{rose1992}
Rose, G. (1992). \textit{The Strategy of Preventive Medicine}. Oxford: Oxford University Press.

\bibitem[Slovic et al.(2004)]{slovic2004}
Slovic, P., Finucane, M.~L., Peters, E., \& MacGregor, D.~G. (2004). Risk as analysis and risk as feelings: some thoughts about affect, reason, risk, and rationality. \textit{Risk Analysis}, 24(2), 311--322.

\bibitem[Small et al.(2007)]{small2007}
Small, D.~A., Loewenstein, G., \& Slovic, P. (2007). Sympathy and callousness: the impact of deliberative thought on donations to identifiable and statistical victims. \textit{Organizational Behavior and Human Decision Processes}, 102(2), 143--153.

\bibitem[Smith(2010)]{smith2010}
Smith, J.~Q. (2010). \textit{Bayesian Decision Analysis: Principles and Practice}. Cambridge: Cambridge University Press.

\bibitem[Tan et al.(2022)]{tan2022}
Tan, W., Colagiuri, B., \& Barnes, K. (2022). Factors moderating the link between personal recounts of COVID-19 vaccine side effects viewed on social media and viewer postvaccination experience. \textit{Vaccines}, 10(10), 1611.

\bibitem[Watson et al.(2022)]{watson2022}
Watson, O.~J., Barnsley, G., Toor, J., Hogan, A.~B., Winskill, P., \& Ghani, A.~C. (2022). Global impact of the first year of COVID-19 vaccination: a mathematical modelling study. \textit{The Lancet Infectious Diseases}, 22(9), 1293--1302.

\end{thebibliography}

\end{document}

